\documentclass[12pt]{amsart}

\usepackage[paper=letterpaper,margin=2cm]{geometry}

% dark mode: \usepackage{xcolor} \pagecolor{black} \color{white}

\parindent=0em
\parskip=1em

\title{Raw computational results}
\author{Jack Schmidt}

\begin{document}

%\appendix
\section{Computed table%: don't include yet
} \label{Appendix}

\newcommand\AGL[2]{\operatorname{AGL}_{#1}\left(#2\right)}
%\newcommand\PSL[2]{\operatorname{PSL}_{#1}\left(#2\right)}
\newcommand\PGL[2]{\operatorname{PGL}_{#1}\left(#2\right)}
\newcommand\PSU[2]{\operatorname{PSU}_{#1}\left(#2^2/#2\right)}
\newcommand\PGU[2]{\operatorname{PGU}_{#1}\left(#2\right)}
\newcommand\PGammaL[2]{\operatorname{P\Gamma{}L}_{#1}\left(#2\right)}
\newcommand\PSigmaL[2]{\operatorname{P\Sigma{}L}_{#1}\left(#2\right)}
\newcommand\PSp[2]{\operatorname{PSp}_{#1}\left(#2\right)}
\newcommand\Alt[1]{\operatorname{A}_{#1}}
\newcommand\Sym[1]{\operatorname{S}_{#1}}
\newcommand\Mathieu[1]{\operatorname{M}_{#1}}
\newcommand\Janko[1]{\operatorname{J}_{#1}}
\newcommand\Sz[1]{\operatorname{Sz}\left(#1\right)}
\newcommand\Aut[1]{\operatorname{Aut}\left(#1\right)}
\newcommand\PrimitiveGroup[4]{\operatorname{Prim}(#1,#2)\cong#3}
\newcommand\PrimitiveGroupShort[4]{\operatorname{Prim}(#1,#2)}
\newcommand\StructureDescription[2]{#1}
\newcommand\WreathProduct[2]{#1 \wr #2}

%The solvablizer of an element $x$ of $G$ is the subset $$\Solv_G(x) = \{ y
%\in G : \langle x,y \rangle \text{ is solvable} \}.$$ If $R(G)$ is a normal
%solvable group of $G$, then $R(G)$ is contained in $\Solv_G(x)$ and $R(G) \Solv_G(x) /
%R(G) = \Solv_{G/R(G)}(xR(R))$. The solvablizer covering number of a group is the minimum
%size of a set $\mathcal{X}$ such that $G = \cup_{x\in\mathcal{X}} \Solv_G(x)$.

We analyze the solvabilizer covering number and the involutionary solvabilizer covering number of every Fitting-free group $G$ with an order less than $\mid M_{22}\mid$. The partial table below showcases our computed results, and we're open to sharing the full table with interested parties. Notably, for direct products, we utilize Theorems~\ref{Car} and \ref{Inv} to determine these covering numbers. % and an
%indication is made whether we can restrict to $\mathcal{X}$ consisting only
%of involutions. 
\cite{GAP4} is used to construct the maximal solvable subgroups and
thus the solvablizer graph. Then a mixed integer linear programming solver
is used to find coverings of minimal size. %A separate column is included
%for the restricted problem of only using solvablizers of involutions,
%which is often quite effective, for example for $\PGL{2}{q}$.

Certain entries are represented as intervals: $[a, b]$ signifies the discovery of a cover using $b$ non-identity elements (or involutions in that column), while it has been proven that no cover employing fewer than $a$ non-identity elements (or involutions) is feasible. The notation $\infty$ is used to indicate the absence of a cover, a circumstance exclusive to the involution column.

%Don't include in the paper directly, but hopefully this helps with
%conjectures.

There is a unique group $G$ such that $S \times T \leq G \leq A \times B$,
$[A:S] = [B:T] = [G:S\times T] = 2$, and $G \neq S \times B, A \times T$.
This group is denoted $A \Yup B$, indicating a direct product, except where
the top is squished. This is only used when $S,T$ are simple, so they
are identified as the socle of $A$ and $B$. In the table, it's important to note that there exists a variety of nonsolvable groups, and due to the extensive range of these groups, not all notations and classifications are exhaustively explained. Other group names are as found in
GAP, for example in the output of StructureDescription. 

We extend our gratitude to Dr. Yuan Zhou at the University of Kentucky for assistance with the integer programming problems.
\vspace{1cm}

\( \begin{array}{rcrc}
\text{Order} & \text{Name} & \alpha & \alpha_{inv} \\ \hline
60 & A_5\cong\PSL{2}{4} & 3 & 3 \\
120 & S_5\cong\PGL{2}{5} & 5 & 5 \\
168 & \PSL{2}{7} & 5 & \infty \\
336 & \PGL{2}{7} & 7 & 7 \\
360 & A_6\cong\PSL{2}{9} & 9 & 9 \\
504 & \PSL{2}{8} & 7 & 7 \\
660 & \PSL{2}{11} & 15 & \infty \\
720 & \Mathieu{10} & 9 & 9 \\
720 & \PGL{2}{9} & 8 & 8 \\
720 & \Sym{6} & 9 & 9 \\
1092 & \PSL{2}{13} & 13 & 13 \\
1320 & \PGL{2}{11} & 11 & 11 \\
1440 & \PGammaL{2}{9} & 9 & 9 \\
1512 & \PGammaL{2}{8} & 7 & 7 \\
2184 & \PGL{2}{13} & 13 & 13 \\
2448 & \PSL{2}{17} & 17 & 17 \\
2520 & \Alt{7} & 41 & \infty \\
3420 & \PSL{2}{19} & 31 & \infty \\
%3600 & \PSL{2}{4} \times \PSL{2}{4} & 3 & 3 \\
4080 & \PSL{2}{16} & 15 & 15 \\
%4896 & \PGL{2}{17} & 17 & 17 \\
5040 & \Sym{7} & 21 & 21 \\

5616 & \PSL{3}{3} & 25 & \infty \\
6048 & \PSU{3}{3} & 49 & \infty \\

6072 & \PSL{2}{23} & [38,41] & \infty \\

%6840 & \PGL{2}{19} & 19 & 19 \\
%7200 & \PSL{2}{4} \times \PGL{2}{5} & 3 & 3 \\
%7200 & \PrimitiveGroup{60}{2}{\Alt{5}^2.2}{A5 : S5} & 5 & 5 \\
7200 & \WreathProduct{\PSL{2}{4}}{2} & 3 & 3 \\

7800 & \PSL{2}{25} & 25 & 25 \\

7920 & \Mathieu{11} & 45 & \infty \\

8160 & \PSL{2}{16}:2 & 17 & 17 \\
9828 & \PSL{2}{27} & [32,48] & [33,49] \\
%10080 & \PSL{2}{7} \times \PSL{2}{4} & 3 & 3 \\

11232 & \Aut{\PSL{3}{3}} & 21 & 21 \\


%12096 & \Aut{\PSU{3}{3}} & 36 & 36 \\
%12144 & \PGL{2}{23} & 23 & 23 \\
12180 & \PSL{2}{29} & 29 & 29 \\

%14400 & \PGL{2}{5} \times \PGL{2}{5} & 5 & 5 \\
%14400 & (\PSL{2}{4} \times \PSL{2}{4}) : (C2 \times C2) & 5 & 5 \\
%14400 & (\PSL{2}{4} \times \PSL{2}{4}) : C4 & 5 & 5 \\
14880 & \PSL{2}{31} & [48,60] & \infty \\

%15600 & \PGL{2}{25} & 25 & 25 \\
%15600 & \PSigmaL{2}{25} & 25 & 25 \\
15600 & \PSL{2}{25}.2 & 25 & 25 \\
%16320 & \PGammaL{2}{16} & 17 & 17 \\
%19656 & \PGL{2}{27} & 26 & 26 \\
20160 & A_8\cong\PSL{4}{2} & 54 & 54 \\ % g_20160_1
20160 & \PSL{3}{4} & 169 & \infty \\ % g_20160_2
%20160 & \PGL{2}{7} \times \PSL{2}{4} & 3 & 3 \\ % g_20160_3
%20160 & \PSL{2}{7} \times \PGL{2}{5} & 5 & 5 \\ % g_20160_4
%20160 & \PSL{2}{7} : \PGL{2}{5} & 5 & 5 \\ % \TransitiveGroup{40}{11251} g_20160_5
%21600 & \PSL{2}{9} \times \PSL{2}{4} & 3 & 3 \\
%24360 & \PGL{2}{29} & 29 & 29 \\
25308 & \PSL{2}{37} & 37 & 37 \\
25920 & \PSp{4}{3} & 9 & 9 \\

%28224 & \PSL{2}{7} \times \PSL{2}{7} & 5 & \infty \\
%28800 & \WreathProduct{\PGL{2}{5}}{2} & 5 & 5 \\
29120 & \Sz{8} & [106,155] & [106,155] \\

%29484 & \PSigmaL{2}{27} & [32,42] & \infty \\
%29760 & \PGL{2}{31} & 31 & 31 \\
%30240 & \PSL{2}{8} \times \PSL{2}{4} & 3 & 3 \\
31200 & \PGammaL{2}{25} & 25 & 25 \\
32736 & \PSL{2}{32} & 31 & 31 \\
34440 & \PSL{2}{41} & 41 & 41 \\
%39600 & \PSL{2}{11} \times \PSL{2}{4} & 3 & 3 \\
39732 & \PSL{2}{43} & [65,93] & \infty \\
40320 & \Sym{8} & [23,24] & [23,24] \\ % g_40320_1

%40320 & \PSigmaL{3}{4} & 129 & 129 \\ % g_40320_3
%40320 & \PrimitiveGroup{56}{3}{PSL(3, 4).2_1}{PSL(3,4) : C2} & [63,96] & [64,96] \\ % g_40320_2
%40320 & \PrimitiveGroup{56}{4}{PSL(3, 4).2_2}{PSL(3,4) : C2} & [48,58] & [50,58] \\ % g_40320_4
%40320 & \PGL{2}{7} \times \PGL{2}{5} & 5 & 5 \\ % g_40320_5
%43200 & \Mathieu{10} \Yup \PGL{2}{5} & 5 & 5 \\ % g_43200_5
%43200 & \Mathieu{10} \times \PSL{2}{4} & 3 & 3 \\ % g_43200_1
%43200 & \PGL{2}{9} \Yup \PGL{2}{5} & 5 & 5 \\ % g_43200_7
%43200 & \PGL{2}{9} \times \PSL{2}{4} & 3 & 3 \\ % g_43200_3
%43200 & \PSL{2}{9} \times \PGL{2}{5} & 5 & 5 \\ % g_43200_4
%43200 & \PSigmaL{2}{9} \Yup \PGL{2}{5} & 5 & 5 \\ % g_43200_6
%43200 & \PSigmaL{2}{9} \times \PSL{2}{4} & 3 & 3 \\ % g_43200_2
\end{array} \)

\(
\begin{array}{rcrc}
\text{Order} & \text{Name} & \alpha & \alpha_{inv} \\ \hline
50616 & \PGL{2}{37} & 37 & 37 \\
%51840 & \Aut{PSp{4}{3}} & 36 & 36 \\
51888 & \PSL{2}{47} & [70,107] & \infty \\
%56448 & \PSL{2}{7} \times \PGL{2}{7} & 5 & 7 \\ % g_56448_1

56448 & \PSL{2}{7} : \PGL{2}{7} & 9 & 22 \\ % g_56448_2
%56448 & \PSL{2}{7} \wr 2 & 5 & 14 \\ % g_56448_3
58800 & \PSL{2}{49} & 49 & 49 \\
%58968 & \PGammaL{2}{27} & 27 & 27 \\
%60480 & \PGL{3}{4} & [70,112] & \infty \\
%60480 & \PSL{2}{8} \times \PGL{2}{5} & 5 & 5 \\
%60480 & \PSL{2}{9} \times \PSL{2}{7} & 5 & 9 \\

62400 & \PSU{3}{4} & 160 & \infty \\
%65520 & \PSL{2}{13} \times \PSL{2}{4} & 3 & 3 \\
%68880 & \PGL{2}{41} & 41 & 41 \\
74412 & \PSL{2}{53} & 53 & 53 \\

%79200 & \PGL{2}{11} \Yup \PGL{2}{5} & 5 & 5 \\ % g_79200_3
%79200 & \PGL{2}{11} \times \PSL{2}{4} & 3 & 3 \\ % g_79200_1
%79200 & \PSL{2}{11} \times \PGL{2}{5} & 5 & 5 \\ % g_79200_2
%79464 & \PGL{2}{43} & 43 & 43 \\
%80640 & \PSL{3}{4} : (2 \times 2) & [58,96] & [62,96] \\
%84672 & \PSL{2}{7} \times \PSL{2}{8} & 5 & 7 \\
%86400 & \Mathieu{10} \times \PGL{2}{5} & 5 & 5 \\
%86400 & \PGL{2}{9} \times \PGL{2}{5} & 5 & 5 \\
%86400 & \Sym{6} \times \PGL{2}{5} & 5 & 5 \\
%86400 & ((A6 : C2) : C2) \times \PSL{2}{4} & 3 & 3 \\
86400 & A5 : ((A6 : C2) : C2) & 5 & 5 \\
%86400 & A5 : ((A6 : C2) : C2) & 5 & 5 \\
%86400 & A5 : ((A6 . C2) : C2) & 5 & 5 \\
%87360 & \Aut{\Sz{8}} & [97,162] & [101,162] \\
%90720 & \PGammaL{2}{8} \times \PSL{2}{4} & 3 & 3 \\
95040 & \Mathieu{12} & [277,327] & \infty \\
102660 & \PSL{2}{59} & [86,143] & \infty \\
103776 & \PGL{2}{47} & 47 & 47 \\
%110880 & \PSL{2}{7} \times \PSL{2}{11} & 5 & \infty \\

%112896 & \PGL{2}{7} \times \PGL{2}{7} & 7 & 7 \\
%112896 & \PrimitiveGroup{64}{70}{\PSL{2}{7}^2.2^2}{(PSL(3,2) x PSL(3,2)) : (C2 x C2)} & 9 & [16,18] \\
%112896 & \PrimitiveGroup{64}{71}{\PSL{2}{7}^2.4}{(PSL(3,2) x PSL(3,2)) : C4} & 9 & [21,22] \\
113460 & \PSL{2}{61} & 61 & 61 \\
117600 & \PGL{2}{49} & 49 & 49 \\
117600 & \PSigmaL{2}{49} & 49 & 49 \\
%117600 & \PrimitiveGroup{50}{6}{\PSL{2}{49}.2_3}{PSL(2,49) . C2} & 49 & 49 \\
%120960 & \Mathieu{10} \Yup \PGL{2}{7} & 9 & 9 \\
%120960 & \Mathieu{10} \times \PSL{2}{7} & 5 & 9 \\
120960 & \PGL{2}{9} \Yup \PGL{2}{7} & 9 & 9 \\
%120960 & \PGL{2}{9} \times \PSL{2}{7} & 5 & 8 \\
%120960 & \PGammaL{3}{4} & [70,118] & [126,129] \\
%120960 & \PSL{2}{9} \times \PGL{2}{7} & 7 & 7 \\
%120960 & \PSigmaL{2}{9} \Yup \PGL{2}{7} & 9 & 9 \\
%120960 & \PSigmaL{2}{9} \times \PSL{2}{7} & 5 & 9 \\
%120960 & \PrimitiveGroup{105}{4}{\PSL{3}{4}.\Sym{3}}{PSL(3,4) : S3} & [57,72] & [58,76] \\
%120960 & \PrimitiveGroup{105}{5}{\PSL{3}{4}.6}{PSL(3,4) : C6} & [47,64] & [49,60] \\

%124800 & \PrimitiveGroup{65}{4}{\PSU{3}{4}.2}{PSU(3,4) : C2} & 80 & 80 \\
126000 & \PSU{3}{5} & [272,429] & \infty \\
%129600 & \PSL{2}{9} \times \PSL{2}{9} & 9 & 9 \\
131040 & \PGL{2}{13} \Yup \PGL{2}{5} & 5 & 5 \\
%131040 & \PGL{2}{13} \times \PSL{2}{4} & 3 & 3 \\
%131040 & \PSL{2}{13} \times \PGL{2}{5} & 5 & 5 \\
%146880 & \PSL{2}{17} \times \PSL{2}{4} & 3 & 3 \\
%148824 & \PGL{2}{53} & 53 & 53 \\
150348 & \PSL{2}{67} & [98,165] & \infty \\
%151200 & \Alt{7} \times \PSL{2}{4} & 3 & 3 \\
%158400 & \PGL{2}{11} \times \PGL{2}{5} & 5 & 5 \\
%163680 & \PGammaL{2}{32} & 31 & 31 \\
%169344 & \PGL{2}{7} \times \PSL{2}{8} & 7 & 7 \\
%172800 & \PGammaL{2}{9} \times \PGL{2}{5} & 5 & 5 \\
175560 & \Janko{1} & [133,249] & [133,249] \\
178920 & \PSL{2}{71} & [104,181] & \infty \\
181440 & \Alt{9} & [40,95] & [56,110] \\
%181440 & \PGammaL{2}{8} \times \PGL{2}{5} & 5 & 5 \\
%181440 & \PSL{2}{9} \times \PSL{2}{8} & 7 & 7 \\
%183456 & \PSL{2}{7} \times \PSL{2}{13} & 5 & 13 \\
190080 & \Mathieu{12} & [74,161] & [74,161] \\
194472 & \PSL{2}{73} & 73 & 73 \\
%205200 & \PSL{2}{19} \times \PSL{2}{4} & 3 & 3 \\
205320 & \PGL{2}{59} & 59 & 59 \\
%216000 & \PSL{2}{4} \times \PSL{2}{4} \times \PSL{2}{4} & 3 & 3 \\
%221760 & \PGL{2}{11} \Yup \PGL{2}{7} & 9 & [29,38] \\
%221760 & \PGL{2}{7} \times \PSL{2}{11} & 7 & 7 \\
%221760 & \PSL{2}{7} \times \PGL{2}{11} & 5 & 11 \\
%225792 & \WreathProduct{\PGL{2}{7}}{2} & 7 & 7 \\
%226920 & \PGL{2}{61} & [60,61] & 61 \\
%235200 & \PGammaL{2}{49} & 49 & 49 \\
%237600 & \PSL{2}{9} \times \PSL{2}{11} & 9 & 9 \\
%241920 & \Aut{\PSL{3}{4}} & [56,70] & [57,70] \\
%241920 & \Mathieu{10} \times \PGL{2}{7} & 7 & 7 \\
%241920 & \PGL{2}{9} \times \PGL{2}{7} & 7 & 7 \\
%241920 & \PGammaL{2}{9} \times \PSL{2}{7} & 5 & 9 \\
%241920 & \PSigmaL{2}{9} \times \PGL{2}{7} & 7 & 7 \\
%241920 & \StructureDescription{\PSL{3}{2} : ((A6 . C2) : C2)}{\PSL{2}{9} \times \PSL{2}{7}} & 9 & 9 \\
%241920 & \StructureDescription{\PSL{3}{2} : ((A6 : C2) : C2)}{\PSL{2}{9} \times \PSL{2}{7}} & 9 & 9 \\
%241920 & \StructureDescription{\PSL{3}{2} : ((A6 : C2) : C2)}{\PSL{2}{9} \times \PSL{2}{7}} & 9 & 9 \\
%244800 & \PSL{2}{16} \times \PSL{2}{4} & 3 & 3 \\
246480 & \PSL{2}{79} & [116,204] & \infty \\
%249600 & \Aut{\PSU{3}{4}} & 80 & 80 \\
%252000 & \Aut{\PSU{3}{5}} & [153,168] & [156,168] \\
%254016 & \PSL{2}{7} \times \PGammaL{2}{8} & 5 & 7 \\
%254016 & \PSL{2}{8} \times \PSL{2}{8} & 7 & 7 \\
%259200 & \Mathieu{10} \Yup \PGL{2}{9} & 9 & 9 \\
%259200 & \Mathieu{10} \Yup \PSigmaL{2}{9} & 9 & 9 \\
%259200 & \Mathieu{10} \times \PSL{2}{9} & 9 & 9 \\
%259200 & \PGL{2}{9} \Yup \PSigmaL{2}{9} & 9 & 9 \\
%259200 & \PGL{2}{9} \times \PSL{2}{9} & 8 & 8 \\
%259200 & \PSigmaL{2}{9} \times \PSL{2}{9} & 9 & 9 \\
%259200 & \PrimitiveGroup{360}{2}{Alt(6)^2.2}{A6 : S6} & 9 & 9 \\
%259200 & \PrimitiveGroup{360}{3}{Alt(6)^2.2}{A6 : (A6 : C2)} & 9 & 9 \\
%259200 & \PrimitiveGroup{360}{4}{Alt(6)^2.2}{A6 : (A6 . C2)} & 9 & 9 \\
%259200 & \WreathProduct{\PSL{2}{9}}{2} & 9 & 9 \\
%262080 & \PGL{2}{13} \times \PGL{2}{5} & 5 & 5 \\
262080 & \PSL{2}{64} & 63 & 63 \\
265680 & \PSL{2}{81} & 81 & 81 \\
285852 & \PSL{2}{83} & [122,218] & \infty \\
%293760 & \PGL{2}{17} \Yup \PGL{2}{5} & 5 & 5 \\
%293760 & \PGL{2}{17} \times \PSL{2}{4} & 3 & 3 \\
%293760 & \PSL{2}{17} \times \PGL{2}{5} & 5 & 5 \\
300696 & \PGL{2}{67} & 67 & 67 \\
%302400 & \Alt{7} \times \PGL{2}{5} & 5 & 5 \\
302400 & \Sym{7} \Yup \PGL{2}{5} & 5 & 5 \\
%302400 & \Sym{7} \times \PSL{2}{4} & 3 & 3 \\
%332640 & \PSL{2}{8} \times \PSL{2}{11} & 7 & 7 \\
%336960 & \PSL{3}{3} \times \PSL{2}{4} & 3 & 3 \\
352440 & \PSL{2}{89} & 89 & 89 \\
357840 & \PGL{2}{71} & [70,71] & [70,71] \\
%362880 & \Mathieu{10} \times \PSL{2}{8} & 7 & 7 \\
%362880 & \PGL{2}{9} \times \PSL{2}{8} & 7 & 7 \\
%362880 & \PSU{3}{3} \times \PSL{2}{4} & 3 & 3 \\
%362880 & \PSigmaL{2}{9} \times \PSL{2}{8} & 7 & 7 \\

362880 & \Sym{9} & 36 & 36 \\
%364320 & \PSL{2}{23} \times \PSL{2}{4} & 3 & 3 \\
%366912 & \PGL{2}{13} \Yup \PGL{2}{7} & 9 & 13 \\
%366912 & \PGL{2}{7} \times \PSL{2}{13} & 7 & 7 \\
%366912 & \PSL{2}{7} \times \PGL{2}{13} & 5 & 13 \\

372000 & \PSL{3}{5} & [261,266] & \infty \\
378000 & \PGU{3}{5} & [271,465] & \infty \\
388944 & \PGL{2}{73} & 73 & 73 \\
%393120 & \PSL{2}{9} \times \PSL{2}{13} & 9 & 9 \\

%410400 & \PGL{2}{19} \Yup \PGL{2}{5} & 5 & 5 \\
%410400 & \PGL{2}{19} \times \PSL{2}{4} & 3 & 3 \\
%410400 & \PSL{2}{19} \times \PGL{2}{5} & 5 & 5 \\
%411264 & \PSL{2}{7} \times \PSL{2}{17} & 5 & 17 \\
%423360 & \PSL{2}{7} \times \Alt{7} & 5 & \infty \\
%432000_1 & \PSL{2}{4} \times \PGammaL{2}{4} \times \PSL{2}{4} & 3 & 3 \\
%432000_2 & \PrimitiveGroup{60}{2}{Alt(5)^2.2}{A5 : S5} \times \PSL{2}{4} & 3 & 3 \\
%432000 & \PrimitiveGroupShort{60}{2}{Alt(5)^2.2}{A5 : S5} \Yup \PrimitiveGroupShort{60}{2}{Alt(5)^2.2}{A5 : S5} \Yup \PrimitiveGroupShort{60}{2}{Alt(5)^2.2}{A5 : S5} & 5 & 5 \\
%432000 & \WreathProduct{\PSL{2}{4}}{2} \Yup \PGammaL{2}{4} & 3 & 3 \\
%432000_5 & \WreathProduct{\PSL{2}{4}}{2} \times \PSL{2}{4} & 3 & 3 \\
%435600 & \PSL{2}{11} \times \PSL{2}{11} & 15 & \infty \\
443520 & \Mathieu{22} & [2181,2501] & \infty \\
%604800 & J_2 & [361,828] & [361,830] \\
\end{array} \)

\vspace{1cm}


\clearpage

\section{More explicit descriptions}

Here are some more explicit descriptions of the smaller groups

\subsection{Order 60}

$\PSL{2}{5}$ of order $60$ is covered by the solvablizers of \textbf{3} of its
\textbf{involutions}, but cannot by covered by the solvablizers of fewer
than $3$ of its non-identity elements. \\
{\tiny $G = \langle (1,2,3), (3,4,5) \rangle$ is covered by the solvablizers
of $\{ \\
\phantom. \quad (1,5)(3,4), \\
\phantom. \quad (1,5)(2,3), \\
\phantom. \quad (1,5)(2,4) \\
\}$}

The maximal soluble subgroups are:
\begin{itemize}
\item 5 conjugate copies of $A_4$
\item 6 conjugate copies of $D_{10}$
\item 10 conjugate copies of $D_{6}$
\end{itemize}

All involutions are conjugate, and each involution is contained in one
copy of $A_4$, two copies of $D_{10}$, and two copies of $D_{6}$, for
a total coverage of 36 elements.

All elements of order 3 are conjugate, and each such element is contained
in two copies of $A_4$, zero copies of $D_{10}$, and one copy of $D_{6}$,
for a total coverage of 24 elements.

All elements of order 5 are $\Aut{G}$ conjugate, and each such element is
contained in zero copies of $A_4$, one copy of $D_{10}$, and zero copies
of $D_{6}$, for a total coverage of 10 elements.

A proof that three elements are necessary (and they must be involutions)
could be given by considering the cyclic subgroups of order 5. There are
six of them, and an involution covers two, an order 3 element covers zero,
and order 5 element covers one. Hence $2x_2 + 0x_3 + 1x_5 \geq 6$ and we
want to minimize $x_2+x_3+x_5$: thus $(x_2,x_3,x_5)=(3,0,0)$ and we should
use three involutions. We then check that there are three involutions that
don't overlap for their order 5 parts, and that the rest works out as well.
That last part strikes me as a pretty explicit calculation that I don't
see how to generalize to an infinite family.

\subsection{Order 120}

$\PGL{2}{5}$ of order $120$ is covered by the solvablizers of \textbf{5} of its
\textbf{involutions}, but cannot by covered by the solvablizers of fewer
than $5$ of its non-identity elements. \\
{\tiny $G=\langle(1,2),(1,2,3,4,5)\rangle$ is covered by the solvablizers of $\{ \\
\phantom. \quad (4,5), \\
\phantom. \quad (3,5), \\
\phantom. \quad (2,5)(3,4), \\
\phantom. \quad (1,2)(3,4), \\
\phantom. \quad (1,5)(3,4) \\
\}$ and by the solvablizers of $\{
(2,3)(4,5), (1,2)(3,4), (1,3)(2,5), (1,4)(3,5), (1,5)(2,4)
\}$ which come form a single conjugacy class of involutions}

The maximal soluble subgroups are: \begin{itemize}
\item 5 conjugate copies of $S_4$
\item 6 conjugate copies of $\AGL{1}{5}$
\item 10 conjugate copies of $D_{12}$
\end{itemize}

There are two conjugacy classes of involutions.

The non-central class has a centralizer of order 12 and 10 conjugates. It
is not the square of an element of $G$. In the standard permutation
representation of $\Sym{5}$, these are the 2-cycles. Each element in
this class is contained in 3 copies of $S_4$, 0 copies of $\AGL{1}{5}$,
4 copies of $D_{12}$. Total coverage of 72.

The central class has a centralizer of order 8 (so it is central in every
Sylow 2-subgroup that contains it) and has 15 members. It is the square of
an element of $G$.  In the standard permutation representation of $\Sym{5}$,
these are the products of 2 disjoint 2-cycles.  Each element in this class
is contained in one copy of $S_4$, 2 copies of $\AGL{1}{5}$, and 2 copies
of $D_{12}$. Total coverage of 72.

All elements of order 3 are conjugate. Each is contained in two copies
of $S_4$, zero copies of $\AGL{1}{5}$, and one copy of $D_{12}$, for a
total coverage of 48.

All elements of order 5 are conjugate. Each is contained in zero copies
of $S_4$, one copy of $\AGL{1}{5}$, and zero copies of $D_{12}$, for
a total coverage of 20.

Again we can consider the cyclic subgroups of order $5$. We need
at least central involutions (or six elements of order five) to
cover them. Indeed, three central involutions is enough to cover them,
and then there are many choices of the next two involutions to cover
the rest of the group. TODO: what parts were missing?

Idea: See if this pattern holds for all $\PSL{2}{p}$ and $\PGL{2}{p}$
for odd prime $p$: consider the Sylow $p$-subgroups. They normalize
nothing, yet are huge, so in the solvablizer graph they are
only connected to elements of their normalizer. Elements of their
normalizer can only normalizer so many Sylow $p$-subgroups (I think
two opposite parabolic subgroups is the idea, and at this low rank,
that is two Sylow $p$-subgroups). Since there are $p+1$ Sylow $p$-subgroups
in $\PSL{2}{p}$ and $\PGL{2}{p}$, we need at least $(p+1)/2$ non-identity
elements of their normalizers to cover them. Hence $\alpha = 3$ is only
for $(p+1)/2 \leq 3 \implies p \leq 5$.

\subsection{Order 168}

$\PSL{2}{7}$ of order $168$ is covered by the solvablizers of \textbf{5}
of its \textbf{non-identity elements}, but cannot by covered by the
solvablizers of fewer than $5$ of its non-identity elements, and cannot
be covered by the solvablizers of its involutions (even an unlimited amount) \\
{\tiny $\{ \\
\phantom. \quad (2,3,8)(4,6,7), \\
\phantom. \quad (1,8)(2,7)(3,5)(4,6),  \\
\phantom. \quad (1,3,5)(2,8,4), \\
\phantom. \quad (1,4,5)(3,7,6), \\
\phantom. \quad (1,6,8)(2,5,7) \\
\}$ }

$\PGL{2}{7}$ of order $336$ is covered by the solvablizers of \textbf{7}
its \textbf{involutions}, but cannot by covered by the solvablizers of
fewer than $7$ of its non-identity elements. \\
{\tiny $\{ \\
\phantom. \quad (2,6)(3,8)(5,7), \\
\phantom. \quad (1,4)(2,5)(3,8)(6,7), \\
\phantom. \quad (1,4)(3,5)(7,8), \\
\phantom. \quad (1,4)(2,3)(5,7)(6,8), \\
\phantom. \quad (1,4)(2,6)(3,7)(5,8), \\
\phantom. \quad (1,4)(2,7)(5,6), \\
\phantom. \quad (1,4)(2,8)(3,6) \\
\}$ }

$\PSL{2}{9}$ of order $360$ is covered by the solvablizers of \textbf{9} of its
\textbf{involutions}, but cannot by covered by the solvablizers of fewer
than $9$ of its non-identity elements. \\
{ \tiny $\{ \\
\phantom. \quad (2,3)(4,5), \\
\phantom. \quad (2,4)(3,5), \\
\phantom. \quad (2,5)(3,4), \\
\phantom. \quad (1,2)(5,6), \\
\phantom. \quad (1,5)(2,6), \\
\phantom. \quad (1,3)(4,6), \\
\phantom. \quad (1,4)(3,6), \\
\phantom. \quad (1,6)(3,4), \\
\phantom. \quad (1,6)(2,5) \\
\}$ }

$\PSL{2}{8}$ of order $504$ is covered by the solvablizers of \textbf{7}
of its \textbf{involutions}, but cannot by covered by the solvablizers of fewer
than $7$ of its non-identity elements. \\
{\tiny $\{
\phantom. \quad (2,5)(3,9)(4,7)(6,8) \\
\phantom. \quad (1,8)(2,9)(3,6)(4,7) \\
\phantom. \quad (1,2)(3,8)(4,7)(5,6) \\
\phantom. \quad (1,5)(2,8)(4,7)(6,9) \\
\phantom. \quad (1,9)(2,3)(4,7)(5,8) \\
\phantom. \quad (1,3)(2,6)(4,7)(5,9) \\
\phantom. \quad (1,6)(3,5)(4,7)(8,9 \\
\}$ }

$\PSL{2}{11}$ of order $660$ is covered by the solvablizers of \textbf{15}
of its \textbf{non-identity elements}, but cannot be covered by the
solvablizers of fewer than 15 of its non-identity elements, and cannot
be covered by the solvablizers of involutions (even an unlimited amount). \\
{\tiny $\{ \\
\phantom. \quad (2,5,7,6,12)(3,10,11,4,8), \\
\phantom. \quad (1,2)(3,4)(5,12)(6,11)(7,10)(8,9), \\
\phantom. \quad (1,4)(2,9)(3,7)(5,12)(6,10)(8,11), \\
\phantom. \quad (1,4)(2,6)(3,11)(5,10)(7,9)(8,12), \\
\phantom. \quad (1,4,9,11,8)(2,3,10,12,7), \\
\phantom. \quad (1,6,11,8,7)(3,9,10,12,5), \\
\phantom. \quad (1,5,7,6,9)(2,10,3,12,4), \\
\phantom. \quad (1,4,6,5,8)(2,9,12,11,3), \\
\phantom. \quad (1,10)(2,9)(3,4)(5,6)(7,11)(8,12), \\
\phantom. \quad (1,5)(2,6)(3,4)(7,8)(9,11)(10,12), \\
\phantom. \quad (1,5,8,6,11)(2,4,9,7,10), \\
\phantom. \quad (1,6)(2,10)(3,8)(4,11)(5,12)(7,9), \\
\phantom. \quad (1,10)(2,12)(3,5)(4,8)(6,11)(7,9), \\
\phantom. \quad (1,10)(2,5)(3,9)(4,7)(6,12)(8,11), \\
\phantom. \quad (1,12)(2,6)(3,5)(4,9)(7,10)(8,11) \\
\}$ }

$\Mathieu{10}$ of order $720$ is covered by the solvablizers of \textbf{9}
of its \textbf{involutions}, but cannot be covered by solvablizers of
fewer than 9 of its non-identity elements. \\
{\tiny $\{ \\
\phantom. \quad (2,4)(3,5)(6,8)(9,10), \\
\phantom. \quad (1,2)(3,6)(5,9)(8,10), \\
\phantom. \quad (1,9)(2,3)(4,8)(5,6), \\
\phantom. \quad (1,3)(2,8)(4,9)(6,10), \\
\phantom. \quad (1,4)(3,10)(5,8)(6,9), \\
\phantom. \quad (1,5)(2,10)(4,6)(8,9), \\
\phantom. \quad (1,10)(2,6)(3,8)(4,5), \\
\phantom. \quad (1,6)(2,5)(3,9)(4,10), \\
\phantom. \quad (1,8)(2,9)(3,4)(5,10) \\
\}$ }

$\Sym{6}$ of order $720$ is covered by the solvablizers of \textbf{9} of its
\textbf{involutions}, but cannot be covered by solvablizers of fewer than
9 of its non-identity elements. \\
{\tiny $\{ \\
\phantom. \quad (3,6)(4,5), \\
\phantom. \quad (1,5)(2,4), \\
\phantom. \quad (1,5)(3,6), \\
\phantom. \quad (2,4)(3,6), \\
\phantom. \quad (2,5)(3,6), \\
\phantom. \quad (1,2)(4,5), \\
\phantom. \quad (1,4)(3,6), \\
\phantom. \quad (1,2)(3,6), \\
\phantom. \quad (1,4)(2,5) \\
\}$ }

$\PGL{2}{9}$ of order $720$ is covered by the solvablizers of \textbf{8}
of its \textbf{involutions}, but cannot be covered by solvablizers of
fewer than 8 of its non-identity elements. \\
{\tiny $\{ \\
\phantom. \quad (2,3)(5,7)(6,8)(9,10), \\
\phantom. \quad (1,2)(3,4)(5,7)(8,10), \\
\phantom. \quad (1,8)(2,9)(3,6)(4,10)(5,7), \\
\phantom. \quad (1,3)(2,4)(5,7)(6,9), \\
\phantom. \quad (1,10)(2,6)(3,9)(4,8)(5,7), \\
\phantom. \quad (1,4)(5,7)(6,10)(8,9), \\
\phantom. \quad (1,9)(2,10)(3,8)(4,6)(5,7), \\
\phantom. \quad (1,6)(2,8)(3,10)(4,9)(5,7) \\
\}$} \\

$\PSL{2}{13}$ of order $1092$ is covered by the solvablizers of \textbf{13}
of its \textbf{involutions}, but cannot be covered by solvablizers of
fewer than 13 of its non-identity elements. \\
{\tiny $\{$ \\
\phantom.\quad $(1,6)(2,13)(5,9)(7,8)(10,14)(11,12),$ \\
\phantom.\quad $(1,10)(2,11)(4,13)(5,8)(7,12)(9,14),$ \\
\phantom.\quad $(2,4)(5,10)(6,14)(7,9)(8,11)(12,13),$ \\
\phantom.\quad $(1,9)(2,7)(4,11)(5,12)(6,10)(8,14),$ \\
\phantom.\quad $(1,14)(4,12)(5,6)(7,10)(8,13)(9,11),$ \\
\phantom.\quad $(1,2)(4,14)(5,13)(6,12)(7,11)(8,10),$ \\
\phantom.\quad $(1,11)(2,9)(4,10)(6,13)(7,14)(8,12),$ \\
\phantom.\quad $(1,7)(2,8)(4,9)(5,14)(6,11)(10,13),$ \\
\phantom.\quad $(1,8)(2,5)(4,7)(6,9)(11,13)(12,14),$ \\
\phantom.\quad $(1,13)(2,10)(4,6)(5,7)(9,12)(11,14),$ \\
\phantom.\quad $(1,4)(2,6)(5,11)(8,9)(10,12)(13,14),$ \\
\phantom.\quad $(1,5)(2,12)(4,8)(6,7)(9,13)(10,11),$ \\
\phantom.\quad $(1,12)(2,14)(4,5)(6,8)(7,13)(9,10)$ \\
$\}$}

$\PGammaL{2}{8}$ of order $1512$ is covered by the solvablizers of
\textbf{7} of its \textbf{involutions}, but cannot be covered by
solvablizers of fewer than 7 of its non-identity elements. \\
{\tiny $\{ \\
\phantom. \quad (2,8)(3,7)(4,9)(5,6) \\
\phantom. \quad (1,2)(3,7)(4,6)(8,9) \\
\phantom. \quad (1,4)(2,5)(3,7)(6,9) \\
\phantom. \quad (1,5)(2,6)(3,7)(4,8) \\
\phantom. \quad (1,6)(2,9)(3,7)(5,8) \\
\phantom. \quad (1,8)(2,4)(3,7)(5,9) \\
\phantom. \quad (1,9)(3,7)(4,5)(6,8) \\
\}$ }

\end{document}
